\documentclass[a4paper,12pt]{article}
\usepackage[utf8]{inputenc}
\usepackage[T1]{fontenc}
\usepackage[ngerman]{babel}
\usepackage{geometry}
\geometry{margin=2.5cm}
\usepackage{hyperref}
\usepackage{booktabs}
\usepackage{amsmath}
\usepackage{parskip}
\setlength{\parskip}{6pt}

\begin{document}

% ============================================================
\section*{U-FNO -- An Enhanced Fourier Neural Operator-Based\\
Deep-Learning Model for Multiphase Flow}
% ============================================================

\textbf{Autoren:} Gege Wen, Zongyi Li, Kamyar Azizzadenesheli,
Anima Anandkumar, Sally M.\ Benson \\
\textbf{Zeitschrift:} \textit{Advances in Water Resources}, 163 (2022), 104180 \\
\textbf{DOI:} \href{https://doi.org/10.1016/j.advwatres.2022.104180}{10.1016/j.advwatres.2022.104180}

% ============================================================
\subsection*{Motivation und Problemstellung}
% ============================================================

Numerische Simulation von Mehrphasenströmungen in porösen Medien ist für
viele geowissenschaftliche Anwendungen unerlässlich -- darunter CO\textsubscript{2}-Speicherung,
Grundwasserströmungen und die Extraktion von Erdöl und Erdgas. Diese Simulationen
erfordern eine feine räumliche und zeitliche Diskretisierung und sind
entsprechend rechenintensiv. Im Kontext probabilistischer Analysen und inverser
Modellierung werden häufig tausende Vorwärtssimulationen benötigt, was
traditionelle Methoden praktisch unerschwinglich macht.

Bestehende datengetriebene Ansätze lassen sich in zwei Kategorien einteilen:
(1) CNN-basierte finite-dimensionale Operatoren, die raumgebundene
Euklidische Abbildungen aus Simulationsdaten erlernen, und (2)
physiksinformierte Ansätze, die Lösungsfunktionen durch neuronale Netze
parametrisieren. Beide Kategorien haben Schwächen: CNN-basierte Modelle
neigen zu Overfitting und benötigen sehr große Trainingsdatensätze, während
physiksinformierte Ansätze für jeden neuen Parametersatz erneut trainiert
werden müssen.

Der Fourier Neural Operator (FNO) von Li et al.\ (2020) adressiert diese
Nachteile durch das Erlernen von Abbildungen in unendlich-dimensionalen
Funktionenräumen via Fast Fourier Transform. FNO-Modelle werden nur einmal
trainiert, sind gitterunabhängig, und weisen eine hervorragende
Generalisierungsfähigkeit auf. Dennoch hat der Standard-FNO eine vergleichsweise
niedrigere Trainingsgenauigkeit als CNN-basierte Modelle, da die Beschränkung
auf wenige Fourier-Moden hochfrequente, lokale Strukturen nicht vollständig
erfassen kann.

% ============================================================
\subsection*{Architektur: Der U-FNO-Ansatz}
% ============================================================

U-FNO kombiniert die Stärken des FNO mit den lokalen Merkmalsextraktions\-fähigkeiten
eines U-Net und adressiert damit die komplementären Schwächen beider Architekturen.
Die Architektur besteht aus drei Schritten:

\begin{enumerate}
    \item \textbf{Lifting:} Die Eingabe $a(x)$ wird über ein vollständig verbundenes
    Netzwerk $P$ auf einen höherdimensionalen Repräsentationsraum angehoben.

    \item \textbf{Iterative Schichten:} Es werden zunächst $L$ Standard-Fourier-Schichten
    angewendet, gefolgt von $M$ neuartigen \emph{U-Fourier-Schichten}. Jede
    U-Fourier-Schicht berechnet:
    \[
        v^{m}_{k+1}(x) := \sigma\!\left(
            \mathcal{K}(v^{m}_{k})(x)
            + \mathcal{U}(v^{m}_{k})(x)
            + W(v^{m}_{k}(x))
        \right),
    \]
    wobei $\mathcal{K}$ der spektrale Kernintegral-Operator (Fourier-Pfad),
    $\mathcal{U}$ ein zweistufiger U-Net-Operator (lokaler CNN-Pfad) und $W$
    eine lineare Biaskorrektur ist.

    \item \textbf{Projektion:} Die gelernte Repräsentation wird über ein weiteres
    vollständiges Netzwerk $Q$ auf den Ausgaberaum projiziert.
\end{enumerate}

Die Autoren zeigen, dass eine hälftige Aufteilung -- gleich viele Fourier- und
U-Fourier-Schichten -- die beste Leistung erzielt. Das resultierende Modell
hat ca.\ 33 Millionen Parameter und ist damit in derselben Größenordnung wie
Standard-FNO und CNN-Benchmark.

% ============================================================
\subsection*{Anwendungsfall und Experimente}
% ============================================================

Das Modell wird auf ein hochkomplexes CO\textsubscript{2}-Wasser-Mehrphasenströmungsproblem
im Kontext geologischer CO\textsubscript{2}-Speicherung angewendet. Simuliert werden
30 Jahre CO\textsubscript{2}-Injektion in heterogene, anisotrope Reservoirs mit weiten
Parameterräumen (u.a.\ Permeabilität, Porosität, Injektionsrate, Reservoirtemperatur).
Die Simulationsdaten werden mit dem kommerziellen Simulator ECLIPSE (e300) generiert.
Der Datensatz umfasst 5\,500 Fälle (4\,500 Training, 500 Validierung, 500 Test).

Als Verlustfunktion wird ein relativer $\ell_p$-Verlust auf das Ausgangsfeld
und seinen Gradienten in radialer Richtung eingesetzt:
\[
    \mathcal{L}(y, \hat{y}) =
    \frac{\|y - \hat{y}\|_p}{\|y\|_p}
    + \beta\,\frac{\|\partial y/\partial r - \partial\hat{y}/\partial r\|_p}{\|\partial y/\partial r\|_p}.
\]
Der Gradiententerm verbessert insbesondere die Vorhersage scharfer Phasengrenzen
(\emph{fronts}) -- ähnlich der Idee, durch physikalische Zusatzterme im Verlust die
Qualität der Lösung in relevanten Bereichen zu steigern.

% ============================================================
\subsection*{Ergebnisse}
% ============================================================

U-FNO übertrifft konsequent alle Vergleichsarchitekturen -- Standard-FNO,
Conv-FNO und CNN-Benchmark -- sowohl für Gassättigung als auch für
Druckaufbau:

\begin{itemize}
    \item \textbf{Gassättigung:} Mittlerer absoluter Fehler im Plume
    (MPE) von 1{,}61\,\% gegenüber 2{,}99\,\% beim CNN ($-46\,\%$);
    $R^2_\mathrm{plume} = 0{,}981$ gegenüber $0{,}955$ beim CNN.

    \item \textbf{Druckaufbau:} Mittlerer relativer Fehler (MRE) von
    0{,}68\,\% gegenüber 0{,}89\,\% beim CNN ($-24\,\%$);
    $R^2 = 0{,}992$ gegenüber $0{,}987$ beim CNN.

    \item \textbf{Dateneffizienz:} U-FNO benötigt nur ein Drittel der
    Trainingsdaten des CNN, um die gleiche Genauigkeit zu erzielen.

    \item \textbf{Rechengeschwindigkeit:} Das trainierte Modell ist
    $6 \times 10^4$-mal schneller als der traditionelle Simulator.

    \item \textbf{Scharfe Phasengrenzen:} U-FNO ist 2{,}7-mal genauer
    als CNN bei der Vorhersage der Gassättigungs-Front und 1{,}8-mal
    genauer bei der Druckaufbau-Front.

    \item \textbf{Heterogene Formationen:} In stark heterogenen Permeabilitätsfeldern
    ist U-FNO 1{,}7-mal genauer als CNN.
\end{itemize}

Bemerkenswert ist weiterhin, dass alle FNO-basierten Modelle ein geringeres
Overfitting aufweisen als der CNN-Benchmark, was die strukturelle
Regularisierungswirkung der Fourier-Schichten bestätigt.

% ============================================================
\subsection*{Bezug zur Masterarbeit}
% ============================================================

Die Arbeit von Wen et al.\ ist aus mehreren Gründen direkt relevant für die
vorliegende Masterarbeit.

\textbf{Architektonische Motivation.}
U-FNO kombiniert exakt die beiden Architekturen, die in dieser Masterarbeit
separat verglichen werden: U-Net und FNO. Die zentrale Erkenntnis der
Arbeit -- dass globale spektrale Merkmalsextraktion (Fourier-Pfad) und lokale
räumliche Verarbeitung (U-Net-Pfad) komplementäre Stärken besitzen -- liefert
eine theoretische Grundlage für die in dieser Masterarbeit durchgeführten
Experimente. Während Wen et al.\ eine hybride Architektur vorschlagen, ist
das primäre Ziel dieser Masterarbeit, die komplementären Eigenschaften beider
Paradigmen durch einen systematischen, kontrollierten Vergleich auf identischen
Datensätzen zu quantifizieren.

\textbf{Komplementäre Stärken der Architekturen.}
U-FNO demonstriert, dass FNO-basierte Modelle besonders bei der Erfassung
großräumiger, elliptischer Strukturen stark sind -- eine Eigenschaft, die
auch im Stokes-Strömungsregime der Masterarbeit von Bedeutung ist, wo lokale
Störungen durch Kristalle sich global im Strömungsfeld auswirken. Die
U-Net-Komponente verbessert hingegen die Auflösung scharfer Phasengrenzen
und lokaler Phänomene. In der Masterarbeit entsprechen diese Rollen der
Erfassung langreichweitiger hydrodynamischer Wechselwirkungen (FNO-Stärke)
versus der präzisen Auflösung von Grenzschichten an Kristalloberflächen
(U-Net-Stärke).

\textbf{Gitterunabhängigkeit und Skalierbarkeit.}
Wen et al.\ diskutieren explizit den Kompromiss zwischen Gitterunabhängigkeit
(eine Stärke des reinen FNO) und Trainingsgenauigkeit (verbessert durch den
U-Net-Pfad). Für die Masterarbeit ist dieser Kompromiss relevant, da alle
Simulationen auf einem festen $256 \times 256$-Gitter durchgeführt werden,
was den Verzicht auf Gitterunabhängigkeit zugunsten höherer Genauigkeit
rechtfertigt.

\textbf{Verlustfunktionsdesign.}
Die Einbeziehung eines Gradiententerms in der Verlustfunktion von Wen et al.\
ist konzeptuell verwandt mit dem in der Masterarbeit eingesetzten
physiksinformierten Ansatz, der über die Strom\-funktion $\psi$ die
Inkompressibilitätsbedingung strukturell erzwingt und über Randterme die
Genauigkeit an Kristalloberflächen erhöht. Beide Ansätze adressieren das
Problem, dass reine Datenverlustfunktionen physikalisch relevante Strukturen
-- scharfe Fronten bzw.\ Grenzschichten -- systematisch unterschätzen.

\textbf{Datensatz und Generalisierung.}
Die Arbeit demonstriert die hervorragende Dateneffizienz FNO-basierter Modelle:
U-FNO benötigt nur ein Drittel der CNN-Trainingsdaten für vergleichbare
Leistung. In der Masterarbeit, wo Trainingsdaten durch rechenintensive
LaMEM-Simulationen begrenzt sind, ist diese Eigenschaft besonders
wertvoll. Die in dieser Masterarbeit untersuchte Generalisierung von
$N$ auf $N+m$ Kristalle stellt zudem eine konzeptuell anspruchsvollere Form
der Out-of-distribution-Generalisierung dar als jene in Wen et al., da sich
nicht nur die Parameterwerte ändern, sondern die topologische Komplexität
des Problems selbst.

\textbf{Verortung im Forschungsfeld.}
U-FNO belegt, dass hybride Architekturen -- die das Beste aus spektraler und
räumlicher Verarbeitung kombinieren -- state-of-the-art Leistung erbringen.
Diese Erkenntnis legt nahe, dass die Diskussion der Ergebnisse der Masterarbeit
explizit auf mögliche Synergiepotenziale zwischen U-Net und FNO eingehen
sollte: Falls die empirischen Experimente zeigen, dass beide Architekturen
komplementäre Stärken aufweisen -- etwa FNO bei großen Kristallzahlen und
globalem Strömungsmuster, U-Net bei lokalen Grenzschichten -- wäre U-FNO
eine natürliche Richtung für zukünftige Arbeit.

\end{document}